% block: matrix-exponential-heading page:210 type:section
\subsection{\BT{Matrix Exponential Function}{矩阵指数函数}}\label{sec:matrix-exponential-function}

% block: matrix-exponential-series-definition page:210 type:ordinary-paragraph
\BT{As seen in Definition~\ref{def:jordan-matrix-function}, the matrix exponential function is defined by the Jordan canonical form. On the other hand, the following equivalent definition is useful for computing the matrix function:}{如 Definition~\ref{def:jordan-matrix-function} 所示, 矩阵指数函数通过 Jordan 标准形定义. 另一方面, 下面的等价定义有助于计算该矩阵函数:}
\begin{equation}\label{eq:matrix-exponential-series}e^A:=I+A+\frac{1}{2!}A^2+\frac{1}{3!}A^3+\cdots.\end{equation}
\BT{The series converges for any square matrix $A\in\mathbb{C}^{n\times n}$. In this section, numerical algorithms for matrix exponential functions are described.}{该级数对任意方阵 $A\in\mathbb{C}^{n\times n}$ 都收敛. 本节介绍矩阵指数函数的数值算法.}

% block: matrix-exponential-properties-transition page:210 type:ordinary-paragraph
\BT{Fundamental properties of matrix functions are summarized in Theorem~\ref{thm:matrix-exponential-properties}. These properties are easily proved from the definition in \eqref{eq:matrix-exponential-series}.}{矩阵函数的基本性质汇总于 Theorem~\ref{thm:matrix-exponential-properties}. 这些性质容易由 \eqref{eq:matrix-exponential-series} 中的定义证明.}

% block: matrix-exponential-properties-theorem page:210 type:theorem
\Needspace{14\baselineskip}
\begin{Theorem}{\BT{Properties of matrix exponential functions}{矩阵指数函数的性质}}\label{thm:matrix-exponential-properties}
\BT{For $A,B\in\mathbb{C}^{n\times n}$:}{对 $A,B\in\mathbb{C}^{n\times n}$:}
\begin{enumerate}[label=(\arabic*)]
\item\label{part:matrix-exp-zero} \BT{$e^O=I$, where $O$ is the zero matrix;}{$e^O=I$, 其中 $O$ 为零矩阵;}
\item\label{part:matrix-exp-commuting-sum} \BT{$e^Ae^B=e^{A+B}$ if $AB=BA$;}{如果 $AB=BA$, 则 $e^Ae^B=e^{A+B}$;}
\item\label{part:matrix-exp-inverse} $e^Ae^{-A}=I$, $(e^A)^{-1}=e^{-A}$;
\item\label{part:matrix-exp-similarity} \BT{$e^{XAX^{-1}}=Xe^AX^{-1}$, where $X$ is a nonsingular matrix;}{$e^{XAX^{-1}}=Xe^AX^{-1}$, 其中 $X$ 为非奇异矩阵;}
\item\label{part:matrix-exp-adjoint} $e^{A^{\mathrm H}}=(e^A)^{\mathrm H}$.
\end{enumerate}
\end{Theorem}

% block: matrix-exponential-hermitian-unitary page:210 type:ordinary-paragraph
\BT{From \ref{part:matrix-exp-adjoint} in Theorem~\ref{thm:matrix-exponential-properties}, matrix exponential function $e^A$ is Hermitian if $A$ is Hermitian. From \ref{part:matrix-exp-inverse} and \ref{part:matrix-exp-adjoint} in Theorem~\ref{thm:matrix-exponential-properties}, matrix exponential function $e^A$ is unitary, i.e., $(e^A)^{\mathrm H}e^A=e^{-A}e^A=I$ if $A$ is skew-Hermitian $(A^{\mathrm H}=-A)$.}{由 Theorem~\ref{thm:matrix-exponential-properties} 的第 \ref{part:matrix-exp-adjoint} 项, 如果 $A$ 是 Hermitian 矩阵, 则矩阵指数函数 $e^A$ 也是 Hermitian 矩阵. 由 Theorem~\ref{thm:matrix-exponential-properties} 的第 \ref{part:matrix-exp-inverse} 和 \ref{part:matrix-exp-adjoint} 项, 如果 $A$ 是斜 Hermitian 矩阵 $(A^{\mathrm H}=-A)$, 则矩阵指数函数 $e^A$ 是酉矩阵, 即 $(e^A)^{\mathrm H}e^A=e^{-A}e^A=I$.}

% block: matrix-exponential-scalar-multiples page:211 type:ordinary-paragraph
\BT{From \ref{part:matrix-exp-commuting-sum} in Theorem~\ref{thm:matrix-exponential-properties} with $A=aC$ and $B=bC$ for scalar values $a$ and $b$, we have $e^{aC}e^{bC}=e^{(a+b)C}$. Here, if $a=1$ and $b=-1$, then we have $e^Ce^{-C}=e^{0C}=e^O=I$ that corresponds to \ref{part:matrix-exp-inverse} in Theorem~\ref{thm:matrix-exponential-properties}.}{在 Theorem~\ref{thm:matrix-exponential-properties} 的第 \ref{part:matrix-exp-commuting-sum} 项中, 对标量 $a$ 与 $b$ 取 $A=aC$ 及 $B=bC$, 得到 $e^{aC}e^{bC}=e^{(a+b)C}$. 这里, 如果 $a=1$ 且 $b=-1$, 则得到 $e^Ce^{-C}=e^{0C}=e^O=I$, 这对应于 Theorem~\ref{thm:matrix-exponential-properties} 的第 \ref{part:matrix-exp-inverse} 项.}

% block: matrix-exponential-algorithms-heading page:211 type:subsection
\subsubsection{\BT{Numerical Algorithms for Matrix Exponential Functions}{矩阵指数函数的数值算法}}\label{subsec:matrix-exponential-algorithms}

% block: matrix-exponential-direct-truncation page:211 type:ordinary-paragraph
\BT{One of the simplest ways to approximately compute matrix exponential functions is to truncate the series in \eqref{eq:matrix-exponential-series}, i.e., $I+A+\frac{1}{2!}A^2+\frac{1}{3!}A^3+\cdots+\frac{1}{m!}A^m$. If all the eigenvalues of matrix $A$ are close to zero, then the truncation error may be small for small $n$. The above approach is, however, inefficient for a matrix $A$ with no such distribution of the eigenvalues of $A$.}{近似计算矩阵指数函数的最简单方法之一是截断 \eqref{eq:matrix-exponential-series} 中的级数, 即\newline{}$I+A+\frac{1}{2!}A^2+\frac{1}{3!}A^3+\cdots+\frac{1}{m!}A^m$.\newline{}如果矩阵 $A$ 的所有特征值都接近于零, 则在 $n$ 较小时截断误差就可能较小. 然而, 对于 $A$ 的特征值不具有这种分布的矩阵 $A$, 上述方法效率不高.}

% block: matrix-exponential-scaled-truncation page:211 type:ordinary-paragraph
\BT{A better and simple way to compute matrix exponential functions is to consider $e^{A/s}$ instead of $e^A$. We see that $(e^{A/s})^s=e^A$ for $s\ge1$, and the all eigenvalues of $A/s$ are $s$ times smaller than the eigenvalues of $A$. Then truncating the following series:}{计算矩阵指数函数的一种更好且简单的方法是考虑 $e^{A/s}$ 而不是 $e^A$. 可以看到, 对 $s\ge1$, 有 $(e^{A/s})^s=e^A$, 且 $A/s$ 的所有特征值都是 $A$ 的特征值的 $s$ 分之一. 于是截断下列级数:}
\begin{equation}\label{eq:matrix-exp-scaled-series}e^{A/s}=I+\frac{1}{1!s}A+\frac{1}{2!s^2}A^2+\frac{1}{3!s^3}A^3+\cdots\end{equation}
\BT{yields}{得到}
\begin{equation}\label{eq:matrix-exp-scaled-truncation}F_{r,s}=I+\frac{1}{1!s}A+\frac{1}{2!s^2}A^2+\frac{1}{3!s^3}A^3+\cdots+\frac{1}{r!s^r}A^r,\end{equation}
\BT{which is an approximation to $(F_{r,s})^s\approx e^A$. Then, the following property holds:}{它给出近似 $(F_{r,s})^s\approx e^A$. 此时有下列性质:}\footnote{\BT{In \BookCite{179}, Theorem~\ref{thm:matrix-exp-truncation-error} is stated on Banach algebra, which is a more general result.}{在 \BookCite{179} 中, Theorem~\ref{thm:matrix-exp-truncation-error} 是在 Banach 代数上陈述的, 是一个更一般的结果.}}

% block: matrix-exp-truncation-error-theorem page:211 type:theorem
\Needspace{14\baselineskip}
\begin{Theorem}{\BookCite{179}}\label{thm:matrix-exp-truncation-error}
\BT{Let $A\in\mathbb{C}^{n\times n}$ and $F_{r,s}$ be the truncation given in \eqref{eq:matrix-exp-scaled-truncation}. Then}{设 $A\in\mathbb{C}^{n\times n}$, $F_{r,s}$ 是 \eqref{eq:matrix-exp-scaled-truncation} 中给出的截断. 则}
\begin{equation}\label{eq:matrix-exp-truncation-error}\|e^A-(F_{r,s})^s\|\le\frac{\|A\|^{r+1}e^{\|A\|}}{s^r(r+1)!},\end{equation}
\BT{where $\|\cdot\|$ is any matrix norm that is submultiplicative.}{其中 $\|\cdot\|$ 是任意具有次乘性的矩阵范数.}\footnotemark
\end{Theorem}
\footnotetext{\BT{See (\ref{cond:matrix-norm-product}) in Section~\ref{subsec:matrix-norm} for the term ``submultiplicative''.}{术语 ``次乘性'' 见 Section~\ref{subsec:matrix-norm} 的 (\ref{cond:matrix-norm-product}).}}

% block: matrix-exp-truncation-parameter-example page:211 type:ordinary-paragraph
\BT{From the right-hand side of Theorem~\ref{thm:matrix-exp-truncation-error}, the truncation error can be estimated by computing $\|A\|$ or its upper bound. For example, we consider the case $s=2^m$. Assume that $\|A\|=10$. Then if we use $r=8$, $m=9$, $s=2^m$, the right-hand side of \eqref{eq:matrix-exp-truncation-error} is about $1.29\times10^{-14}$, which means that it is possible to determine $s$ and $r$ such that the truncation error is less than a given tolerance, if $\|A\|$ is estimated. When $\|A\|$ is a Frobenius norm, then it is easy to compute the norm by the definition of a Frobenius norm. When $\|A\|$ is a matrix 2-norm, the maximum singular value of $A$ is needed, which can also be easily estimated by the Golub--Kahan bidiagonalization process (Algorithm~\ref{alg:golub-kahan}).}{根据 Theorem~\ref{thm:matrix-exp-truncation-error} 的右端, 可以通过计算 $\|A\|$ 或其上界来估计截断误差. 例如, 考虑 $s=2^m$ 的情形. 假设 $\|A\|=10$. 那么, 若取 $r=8$, $m=9$, $s=2^m$, 则 \eqref{eq:matrix-exp-truncation-error} 的右端约为 $1.29\times10^{-14}$. 这意味着, 如果估计出 $\|A\|$, 就可以确定 $s$ 与 $r$, 使截断误差小于给定容差. 当 $\|A\|$ 是 Frobenius 范数时, 容易根据 Frobenius 范数的定义计算该范数. 当 $\|A\|$ 是矩阵 2-范数时, 需要 $A$ 的最大奇异值, 它也容易通过 Golub--Kahan 双对角化过程 (Algorithm~\ref{alg:golub-kahan}) 估计.}

% block: matrix-exp-pade-approximant page:212 type:ordinary-paragraph
\BT{A more cost-efficient method is as follows: instead of computing \eqref{eq:matrix-exp-scaled-truncation}, we use the Padé approximant to $e^{A/s}$ and compute $s$th power of the approximant. The Padé approximant $R_{p,q}(A)$ to matrix exponential function $e^A$ is given below:}{一种计算成本更低的方法如下: 不计算 \eqref{eq:matrix-exp-scaled-truncation}, 而采用 $e^{A/s}$ 的 Padé 逼近, 并计算该逼近的 $s$ 次幂. 矩阵指数函数 $e^A$ 的 Padé 逼近 $R_{p,q}(A)$ 如下:}
\begin{equation}\label{eq:matrix-exp-pade}R_{p,q}(A)=[D_{p,q}(A)]^{-1}N_{p,q}(A),\end{equation}
\BT{where}{其中}
\[\begin{aligned}N_{p,q}(A)&=\sum_{j=0}^p n_jA^j,\quad n_j=\frac{(p+q-j)!}{(p+q)!}\times\frac{p!}{j!(p-j)!},\\D_{p,q}(A)&=\sum_{j=0}^q d_j(-A)^j,\quad d_j=\frac{(p+q-j)!}{(p+q)!}\times\frac{q!}{j!(q-j)!}.\end{aligned}\]

% block: matrix-exp-diagonal-pade-parameters page:212 type:ordinary-paragraph
\BT{When approximating the matrix exponential, it is recommended to use the diagonal Padé approximant, i.e., $p=q$. From $e^A=[e^{A/2^m}]^{2^m}\approx[R_{q,q}(A/2^m)]^{2^m}$, we need to choose two parameters $q$ and $m$. If $\|A\|/2^m\le1/2$, then}{近似矩阵指数时, 建议采用对角 Padé 逼近, 即 $p=q$. 由\newline{}$e^A=[e^{A/2^m}]^{2^m}\approx[R_{q,q}(A/2^m)]^{2^m}$,\newline{}需要选择两个参数 $q$ 与 $m$. 如果 $\|A\|/2^m\le1/2$, 则}
\[[R_{q,q}(A/2^m)]^{2^m}=e^{A+E},\quad\frac{\|E\|}{\|A\|}\le8\left[\frac{\|A\|}{2^m}\right]^{2q}\frac{(q!)^2}{(2q)!(2q+1)!}.\]
\BT{For the details, see [\hyperref[bib:source-133]{\ref*{bib:source-133}}, p.12]. The inequality leads to optimum parameters $q$ and $m$ such that $q+m$ is minimized under the condition that the truncation error holds $\|E\|/\|A\|\le\epsilon$ for a given tolerance $\epsilon$. The optimum parameters for a tolerance $\epsilon=10^{-15}$ with respect to $\|A\|$ are given in Table~\ref{tab:matrix-exp-pade-parameters}. For related studies, see \BookCite{8}, \BookCite{98}.}{详情见 [\hyperref[bib:source-133]{\ref*{bib:source-133}}, p.12]. 该不等式给出最优参数 $q$ 与 $m$, 使得在对给定容差 $\epsilon$ 截断误差满足 $\|E\|/\|A\|\le\epsilon$ 的条件下, $q+m$ 最小. 对容差 $\epsilon=10^{-15}$, 随 $\|A\|$ 变化的最优参数列于 Table~\ref{tab:matrix-exp-pade-parameters}. 相关研究见 \BookCite{8}, \BookCite{98}.}

% block: matrix-exp-schur-reduction page:212 type:ordinary-paragraph
\BT{If Schur decomposition $A=QTQ^{\mathrm H}$ is computed, then it follows from \eqref{eq:matrix-exponential-series} that $e^A=e^{QTQ^{\mathrm H}}=Qe^TQ^{\mathrm H}$. Thus all we have to do is to compute $e^T$ using (diagonal) Padé approximant.}{如果计算出 Schur 分解 $A=QTQ^{\mathrm H}$, 则由 \eqref{eq:matrix-exponential-series} 可得 $e^A=e^{QTQ^{\mathrm H}}=Qe^TQ^{\mathrm H}$. 因此, 只需利用 (对角) Padé 逼近计算 $e^T$.}

% block: matrix-exp-hermitian-eigen-decomposition page:212 type:ordinary-paragraph
\BT{In particular, when $A$ is Hermitian, $T$ becomes a diagonal matrix. Let $\lambda_1,\lambda_2,\ldots,\lambda_n$ be the diagonal elements of $T$. Then the matrix exponential function can be computed by}{特别地, 当 $A$ 是 Hermitian 矩阵时, $T$ 成为对角矩阵. 设 $\lambda_1,\lambda_2,\ldots,\lambda_n$ 是 $T$ 的对角元素. 则矩阵指数函数可由下式计算:}
\begin{equation}\label{eq:matrix-exp-hermitian-diagonal}e^A=Q\begin{bmatrix}e^{\lambda_1}&&\\&\ddots&\\&&e^{\lambda_n}\end{bmatrix}Q^{\mathrm H}.\end{equation}
\BT{Further, if $A$ is diagonalizable, i.e., $A=XDX^{-1}$ for a nonsingular matrix $X$ and a diagonal matrix $D$, then we have $e^A=Xe^DX^{-1}$, where $e^D$ is the same form as in \eqref{eq:matrix-exp-hermitian-diagonal}. On the other hand, if the condition number $X$ is high, this approach can be numerically unstable.}{进一步, 如果 $A$ 可对角化, 即对非奇异矩阵 $X$ 和对角矩阵 $D$ 有 $A=XDX^{-1}$, 则 $e^A=Xe^DX^{-1}$, 其中 $e^D$ 与 \eqref{eq:matrix-exp-hermitian-diagonal} 中的形式相同. 另一方面, 如果 $X$ 的条件数很大, 这种方法在数值上可能不稳定.}

% block: matrix-exp-pade-parameter-table page:212 type:table
\begin{table}[H]\centering
\caption{\BT{Optimum parameters $q$ and $m$ in $R_{q,q}(A/2^m)$ satisfying the error tolerance $\epsilon=10^{-15}$.}{满足误差容差 $\epsilon=10^{-15}$ 的 $R_{q,q}(A/2^m)$ 中的最优参数 $q$ 与 $m$.}}\label{tab:matrix-exp-pade-parameters}
\begin{tabular}{c|cccccccccc}\hline
$\|A\|$&$10^{-3}$&$10^{-2}$&$10^{-1}$&$1$&$10$&$10^2$&$10^3$&$10^4$&$10^5$&$10^6$\\\hline
$q$&3&3&4&6&6&6&6&6&6&6\\
$m$&0&0&0&1&5&8&11&15&18&21\\\hline
\end{tabular}\end{table}

% block: matrix-exponential-vector-heading page:213 type:subsection
\subsubsection{\BT{Multiplication of a Matrix Exponential Function and a Vector}{矩阵指数函数与向量的乘积}}\label{subsec:matrix-exponential-vector}

% block: matrix-exponential-ode-solution page:213 type:ordinary-paragraph
\BT{We consider the following linear ordinary differential equation:}{考虑如下线性常微分方程:}
\begin{equation}\label{eq:matrix-exp-linear-ode}\frac{\mathrm{d}\boldsymbol{x}(t)}{\mathrm{d}t}=A\boldsymbol{x}(t)+\boldsymbol{f}(t),\quad\boldsymbol{x}(0)=\boldsymbol{x}_0,\end{equation}
\BT{where $A\in\mathbb{C}^{n\times n}$, $\boldsymbol{x},\boldsymbol{f}\in\mathbb{C}^n$, $t\in\mathbb{R}$. The solution can be given by using matrix exponential functions as follows:}{其中 $A\in\mathbb{C}^{n\times n}$, $\boldsymbol{x},\boldsymbol{f}\in\mathbb{C}^n$, $t\in\mathbb{R}$. 利用矩阵指数函数可如下给出其解:}
\begin{equation}\label{eq:matrix-exp-variation-of-constants}\boldsymbol{x}(t)=e^{tA}\boldsymbol{x}_0+\int_0^t e^{(t-\tau)A}\boldsymbol{f}(\tau)\,\mathrm{d}\tau.\end{equation}

% block: matrix-exponential-vector-product-motivation page:213 type:ordinary-paragraph
\BT{We see that in \eqref{eq:matrix-exp-variation-of-constants} there are multiplications of matrix exponential functions and vectors. In particular, if $\boldsymbol{f}(t)=\boldsymbol{0}$, then the solution is written as $\boldsymbol{x}(t)=e^{tA}\boldsymbol{x}_0$. This explicit form of the solution means that we know the solution at an arbitrary time $t$ if the multiplications are obtained.}{可以看到, \eqref{eq:matrix-exp-variation-of-constants} 中出现了矩阵指数函数与向量的乘积. 特别地, 如果 $\boldsymbol{f}(t)=\boldsymbol{0}$, 则解可写为 $\boldsymbol{x}(t)=e^{tA}\boldsymbol{x}_0$. 解的这一显式形式意味着, 如果得到这些乘积, 就知道任意时刻 $t$ 的解.}

% block: matrix-exponential-vector-cost page:213 type:ordinary-paragraph
\BT{If matrix exponential functions have already been computed, then all we have to do is to compute the matrix--vector multiplication whose computational cost is of order $n^2$, where $n$ is a matrix of size $n\times n$. On the other hand, computing the cost of a matrix exponential function may be of order $n^3$ even if matrix $A$ is sparse. Thus computing matrix exponential functions is inefficient, when it comes to computing the solution of \eqref{eq:matrix-exp-linear-ode}. Note that the solution only requires the result of the multiplication of a matrix exponential functions and a vector. In such a case, an approximation to the multiplication using Krylov subspace is useful, which is described below.}{如果已经计算出矩阵指数函数, 则只需计算矩阵与向量的乘积, 其计算成本为 $n^2$ 阶, 其中 $n$ 是一个大小为 $n\times n$ 的矩阵. 另一方面, 即使矩阵 $A$ 是稀疏的, 计算矩阵指数函数的成本也可能为 $n^3$ 阶. 因此, 在求解 \eqref{eq:matrix-exp-linear-ode} 时计算矩阵指数函数并不高效. 注意, 该解只需要矩阵指数函数与向量相乘的结果. 在这种情况下, 利用 Krylov 子空间近似该乘积是有用的, 下面予以介绍.}

% block: matrix-exponential-krylov-least-squares page:213 type:ordinary-paragraph
\BT{Since matrix exponential function $e^{tA}$ is defined by the series \eqref{eq:matrix-exponential-series}, it is natural to use matrix polynomials $p_m(A)$ of degree $m$ to approximate $e^{tA}$. The multiplication of a matrix polynomial $p_m(A)$ and a vector belongs to Krylov subspace $\mathcal{K}_m(A,\boldsymbol{x}_0)$. Thus the optimal approximation is given by solving the following least-squares problem:}{由于矩阵指数函数 $e^{tA}$ 由级数 \eqref{eq:matrix-exponential-series} 定义, 自然可以利用 $m$ 次矩阵多项式 $p_m(A)$ 近似 $e^{tA}$. 矩阵多项式 $p_m(A)$ 与向量相乘的结果属于 Krylov 子空间 $\mathcal{K}_m(A,\boldsymbol{x}_0)$. 因此, 最优近似由求解下列最小二乘问题给出:}
\begin{equation}\label{eq:matrix-exp-krylov-least-squares}\min_{\boldsymbol{x}_m\in\mathcal{K}_m(A,\boldsymbol{x}_0)}\|e^{tA}\boldsymbol{x}_0-\boldsymbol{x}_m\|.\end{equation}

% block: matrix-exponential-orthogonal-projection page:213 type:ordinary-paragraph
\BT{Let $V_m$ be a matrix whose columns are orthonormalized basis vectors of $\mathcal{K}_m(A,\boldsymbol{x}_0)$. Then, any vector $\boldsymbol{x}\in\mathcal{K}_m(A,\boldsymbol{x}_0)$ can be written as $\boldsymbol{x}=V_m\boldsymbol{y}$, where $\boldsymbol{y}\in\mathbb{C}^m$. Thus \eqref{eq:matrix-exp-krylov-least-squares} is equivalent to}{设 $V_m$ 是以 $\mathcal{K}_m(A,\boldsymbol{x}_0)$ 的标准正交化基向量为列的矩阵. 则任意向量 $\boldsymbol{x}\in\mathcal{K}_m(A,\boldsymbol{x}_0)$ 均可写为 $\boldsymbol{x}=V_m\boldsymbol{y}$, 其中 $\boldsymbol{y}\in\mathbb{C}^m$. 因此 \eqref{eq:matrix-exp-krylov-least-squares} 等价于}
\[\min_{\boldsymbol{y}\in\mathbb{C}^m}\|e^{tA}\boldsymbol{x}_0-V_m\boldsymbol{y}\|.\]
\BT{The solution is $\boldsymbol{y}=V_m^{\mathrm H}e^{tA}\boldsymbol{x}_0$, and thus the solution of \eqref{eq:matrix-exp-krylov-least-squares} is given by}{其解为 $\boldsymbol{y}=V_m^{\mathrm H}e^{tA}\boldsymbol{x}_0$, 因此 \eqref{eq:matrix-exp-krylov-least-squares} 的解为}
\begin{equation}\label{eq:matrix-exp-krylov-projection}(\boldsymbol{x}(t)\approx)\ \boldsymbol{x}_m=V_mV_m^{\mathrm H}e^{tA}\boldsymbol{x}_0.\end{equation}

% block: matrix-exponential-arnoldi-approximation page:214 type:ordinary-paragraph
\BT{In what follows, we consider using the Arnoldi process in Section~\ref{subsec:arnoldi-process} to compute basis vectors of $\mathcal{K}_m(A,\boldsymbol{x}_0)$ and give an approximate solution of \eqref{eq:matrix-exp-krylov-projection}. From the Arnoldi process in Section~\ref{subsec:arnoldi-process}, it follows from $\boldsymbol{x}_0=\|\boldsymbol{x}_0\|V_m\boldsymbol{e}_1$ that we obtain $\boldsymbol{x}_m=\|\boldsymbol{x}_0\|V_m(V_m^{\mathrm H}e^{tA}V_m)\boldsymbol{e}_1$, where $\boldsymbol{e}_1=[1,0,\ldots,0]^{\top}$. Furthermore, using the relation $V_m^{\mathrm H}AV_m=H_m$ (Hessenberg matrix) in \eqref{eq:arnoldi-projected-matrix}, we consider an approximation $V_m^{\mathrm H}e^{tA}V_m\approx e^{tV_m^{\mathrm H}AV_m}=e^{tH_m}$. Then we have}{以下考虑利用 Section~\ref{subsec:arnoldi-process} 中的 Arnoldi 过程计算 $\mathcal{K}_m(A,\boldsymbol{x}_0)$ 的基向量, 并给出 \eqref{eq:matrix-exp-krylov-projection} 的近似解. 根据 Section~\ref{subsec:arnoldi-process} 中的 Arnoldi 过程, 由 $\boldsymbol{x}_0=\|\boldsymbol{x}_0\|V_m\boldsymbol{e}_1$ 可得\newline{}$\boldsymbol{x}_m=\|\boldsymbol{x}_0\|V_m(V_m^{\mathrm H}e^{tA}V_m)\boldsymbol{e}_1$,\newline{}其中 $\boldsymbol{e}_1=[1,0,\ldots,0]^{\top}$. 进一步, 利用 \eqref{eq:arnoldi-projected-matrix} 中的关系 $V_m^{\mathrm H}AV_m=H_m$ (Hessenberg 矩阵), 考虑近似\newline{}$V_m^{\mathrm H}e^{tA}V_m\approx e^{tV_m^{\mathrm H}AV_m}=e^{tH_m}$.\newline{}于是有}
\begin{equation}\label{eq:matrix-exp-arnoldi-approximation}\boldsymbol{x}_m=\|\boldsymbol{x}_0\|V_m(V_m^{\mathrm H}e^{tA}V_m)\boldsymbol{e}_1\approx\|\boldsymbol{x}_0\|V_me^{tH_m}\boldsymbol{e}_1.\end{equation}

% block: matrix-exponential-projected-cost page:214 type:ordinary-paragraph
\BT{$e^{tH_m}\boldsymbol{e}_1$ is the multiplication of an $m\times m$ small matrix exponential function and a vector, whose computational costs are of $O(m^2)$. This comes from the fact that the multiplication of the Padé approximant (to $e^{tH_m}$) and a vector requires $H_m^{-1}\boldsymbol{v}$, i.e., solving linear systems $H_m\boldsymbol{z}=\boldsymbol{v}$, whose computational cost is of $O(m^2)$.}{表达式 $e^{tH_m}\boldsymbol{e}_1$ 是一个 $m\times m$ 小型矩阵指数函数与向量的乘积, 其计算成本为 $O(m^2)$. 这是因为 ($e^{tH_m}$ 的) Padé 逼近与向量相乘需要计算 $H_m^{-1}\boldsymbol{v}$, 即求解线性方程组 $H_m\boldsymbol{z}=\boldsymbol{v}$, 其计算成本为 $O(m^2)$.}

% block: matrix-exp-hermitian-error-transition page:214 type:ordinary-paragraph
\BT{If matrix $A$ is Hermitian, a suitable choice of $m$ for satisfying a given error tolerance is provided by using Theorem~\ref{thm:matrix-exp-krylov-error}.}{如果矩阵 $A$ 是 Hermitian 矩阵, 则可利用 Theorem~\ref{thm:matrix-exp-krylov-error} 选择适当的 $m$, 以满足给定的误差容差.}

% block: matrix-exp-hermitian-error-theorem page:214 type:theorem
\Needspace{14\baselineskip}
\begin{Theorem}{\BookCite{99}, Theorem 2}\label{thm:matrix-exp-krylov-error}
\BT{Let $A$ be Hermitian and all the eigenvalues lie in the interval $[-4\rho,0]$. Then the error of \eqref{eq:matrix-exp-arnoldi-approximation} is given by}{设 $A$ 是 Hermitian 矩阵, 其所有特征值都位于区间 $[-4\rho,0]$ 中. 则 \eqref{eq:matrix-exp-arnoldi-approximation} 的误差为}
\[\|e^{tA}\boldsymbol{x}_0-\|\boldsymbol{x}_0\|_2V_me^{tH_m}\boldsymbol{e}_1\|\le\begin{cases}c\cdot e^{-\frac{m^2}{5\rho t}}&(\sqrt{4\rho t}\le m\le2\rho t),\\c\cdot\frac{1}{\rho te^{\rho t}}\cdot\left(\frac{e\rho t}{m}\right)^m&(2\rho t\le m),\end{cases}\]
\BT{where $c=10\|\boldsymbol{x}_0\|$.}{其中 $c=10\|\boldsymbol{x}_0\|$.}
\end{Theorem}

% block: matrix-exp-practical-dimension-selection page:214 type:ordinary-paragraph
\BT{In practice, first estimate the minimum eigenvalue $\lambda_{\min}$ of $A$ and set $\rho=-\lambda_{\min}/4$. Then, the dimension of Krylov subspace $m$ is chosen so that the upper bound error in Theorem~\ref{thm:matrix-exp-krylov-error} is less than the given tolerance, e.g., $10^{-12}$. The usefulness of the error bound is shown in \BookCite{139}, together with devising a practical computation avoiding the loss of orthogonality regarding the basis vectors of Krylov subspaces and an efficient method of parallel computation.}{实际计算中, 先估计 $A$ 的最小特征值 $\lambda_{\min}$, 并设置 $\rho=-\lambda_{\min}/4$. 然后选择 Krylov 子空间的维数 $m$, 使得 Theorem~\ref{thm:matrix-exp-krylov-error} 中的误差上界小于给定容差, 例如 $10^{-12}$. 文献 \BookCite{139} 展示了该误差界的有效性, 并设计了避免 Krylov 子空间基向量丧失正交性的实用计算方法和高效的并行计算方法.}
